\documentclass{article}
\usepackage{amsmath,amssymb,bm,booktabs,graphicx}
\begin{document}
\title{Adaptive Kernels for Sample-Efficient Optimization}
\author{A. Author}
\maketitle

\begin{abstract}
Gaussian processes (GPs) are the standard surrogate for Bayesian optimization (BO), but a fixed kernel limits
sample efficiency. We propose AKO, a method that adapts the kernel online. AKO reduces regret by 41\% on
average and outperforms all baselines. We also show that its regret is sublinear.
\end{abstract}

\section{Introduction}
GPs model the objective with a prior over functions. We use expected improvement (EI) to pick the next point.
Our method achieves state-of-the-art results across benchmarks. KRR is a common baseline in this literature.

\section{Method}
Let $f$ be the objective function and let $\mathbf{x}$ denote the input. We observe $n$ noisy evaluations,
where $n$ is the number of observations. The surrogate uses a kernel, where $k$ is the kernel function.
Let $K$ denote the kernel function evaluated on the training inputs.

We write the noise as $\epsilon$ with variance $\sigma^2$. The posterior mean at $\boldsymbol{x}$ is
\begin{equation}
  \mu(\mathbf{x}) = k(\mathbf{x}, X)^\top K^{-1} y, \qquad \varepsilon \sim \mathcal{N}(0, \sigma^2),
\end{equation}
where $n$ is the input dimension. Here $\bm{x}$ is any query point. EI is the expected improvement (EI) over the
incumbent. We treat the kernel hyperparameters as random.

\begin{theorem}\label{thm:regret}
The cumulative regret of AKO is sublinear in the number of evaluations.
\end{theorem}


\section{Experiments}
Table~\ref{tab:main} compares AKO with the baselines on five benchmarks. AKO has the lowest regret on four of them.
On the hardest benchmark AKO is significantly faster than random search.
We also report the sensitivity to the budget, which is flat beyond 50 evaluations.

\begin{table}[t]
\caption{Final regret, lower is better. AKO has the lowest regret on four of five benchmarks.}
\label{tab:main}
\centering
\begin{tabular}{lcc}
\toprule
Benchmark & Random search & AKO \\
\midrule
Branin & 0.42 & 0.11 \\
Hartmann & 1.30 & 0.84 \\
\bottomrule
\end{tabular}
\end{table}

\begin{figure}[t]
\caption{Sensitivity of AKO to the evaluation budget.}
\label{fig:budget}
\end{figure}

\section{Conclusion}
We introduced AKO, see Figure~\ref{fig:missing}. AKO reduces regret by 23\% over the best baseline.
\end{document}
